\chapter{Reference}
\section{Giới thiệu về Ngôn ngữ Truy vấn KBQL}

KBQL (Knowledge Base Query Language) là ngôn ngữ truy vấn chính được sử dụng trong hệ quản trị KBMS. KBQL không chỉ kế thừa các cú pháp SQL tiêu chuẩn để thao tác với dữ liệu mà còn mở rộng các khả năng suy diễn tri thức (Reasoning) dựa trên logic vị từ và tập luật (Rules).

\subsection{Triết lý Thiết kế}

KBQL được thiết kế dựa trên ba trụ cột chính:
1.  \textbf{Sự quen thuộc (Familiarity):} Cú pháp gần gũi với SQL giúp người dùng dễ dàng tiếp cận.
2.  \textbf{Tính tri thức (Knowledge-Driven):} Tích hợp sâu các khái niệm về Concept, Fact và Rule thay vì chỉ dừng lại ở Table/Row.
3.  \textbf{Tự động suy diễn (Automatic Inference):} Kết quả truy vấn có thể được tự động cập nhật hoặc suy luận thông qua bộ máy Suy diễn (Inference Engine) mà người dùng không cần viết logic thủ công.

\subsection{Các Thành phần Chính của KBQL}

| Nhóm Lệnh | Chức năng | Các lệnh tiêu biểu |
| :--- | :--- | :--- |
| \textbf{KDL} (Knowledge Definition) | Định nghĩa cấu trúc tri thức, luật, phân cấp | `CREATE KB`, `CONCEPT`, `RULE`, `HIERARCHY`, `RELATION` |
| \textbf{KML} (Knowledge Manipulation) | Thao tác trên tập các sự kiện (Facts) | `INSERT`, `UPDATE`, `DELETE`, `IMPORT`, `EXPORT` |
| \textbf{KQL} (Knowledge Query) | Truy vấn và yêu cầu suy diễn | `SELECT`, `SOLVE`, `SHOW`, `EXPLAIN`, `DESCRIBE` |
| \textbf{KCL} (Knowledge Control) | Quản lý người dùng và quyền truy cập | `GRANT`, `REVOKE`, `CREATE/ALTER/DROP USER` |
| \textbf{TCL} (Transaction Control) | Quản lý giao dịch và tính toàn vẹn | `BEGIN`, `COMMIT`, `ROLLBACK` |
| \textbf{Admin} (Maintenance) | Bảo trì và tối ưu hóa hệ thống | `MAINTENANCE (VACUUM, REINDEX, CHECK)` |

\subsection{Hỗ trợ Chỉnh sửa (ALTER Support)}

Không phải tất cả các đối tượng trong KBMS đều hỗ trợ lệnh `ALTER`. Dưới đây là bảng tra cứu:

| Đối tượng | Hỗ trợ ALTER | Ghi chú |
| :--- | :--- | :--- |
| \textbf{Concept} | ✅ Có | Hỗ trợ thêm/xóa biến, luật, ràng buộc, quan hệ nội bộ. |
| \textbf{Knowledge Base} | ✅ Có | Hỗ trợ thay đổi mô tả (Description). |
| \textbf{User} | ✅ Có | Hỗ trợ đổi mật khẩu và vai trò quản trị. |
| \textbf{Relation} | ❌ Không | Cần `DROP` và `CREATE` lại. |
| \textbf{Hierarchy} | ❌ Không | Sử dụng `ADD/REMOVE HIERARCHY`. |
| \textbf{Rule} (Toàn cục) | ❌ Không | Cần `DROP` và `CREATE` lại (Khác với Rule nội bộ Concept). |
| \textbf{Operator/Function}| ❌ Không | Cần `DROP` và `CREATE` lại. |

\subsection{Kiểu Dữ liệu Hỗ trợ (Data Types)}

KBQL cung cấp bộ kiểu dữ liệu phong phú để định nghĩa Concept:

| Nhóm | Kiểu dữ liệu | Mô tả |
| :--- | :--- | :--- |
| \textbf{Số học} | `INT`, `BIGINT`, `DECIMAL`, `FLOAT`, `DOUBLE` | Các kiểu số nguyên và số thực. |
| \textbf{Chuỗi} | `VARCHAR(n)`, `CHAR(n)`, `TEXT`, `STRING` | Lưu trữ văn bản (hỗ trợ độ dài tùy chỉnh). |
| \textbf{Logic} | `BOOLEAN` | Giá trị `true` hoặc `false`. |
| \textbf{Thời gian} | `DATE`, `DATETIME`, `TIMESTAMP` | Quản lý thời gian và sự kiện. |
| \textbf{Tri thức} | `OBJECT`, `<ConceptName>` | Tham chiếu đến một đối tượng hoặc một Khái niệm khác. |
| \textbf{Đặc biệt} | `NULL` | Trạng thái rỗng. |

\subsection{Khái niệm Cốt lõi}

\subsubsection{Concept (Khái niệm)}
Thay vì dùng "Table", KBQL sử dụng \textbf{Concept}. Một Concept đại diện cho một thực thể hoặc một lớp đối tượng trong thế giới thực, bao gồm các biến (Variables) định nghĩa thuộc tính của nó.

\subsubsection{Fact (Sự kiện)}
Mỗi bản ghi dữ liệu trong một Concept được coi là một \textbf{Fact}. Tập hợp các Fact tạo nên cơ sở dữ liệu hiện tại.

\subsubsection{Rule (Luật)}
\textbf{Rule} định nghĩa mối quan hệ logic giữa các Fact. Khi dữ liệu (Fact) thay đổi, các Rule liên quan có thể được kích hoạt để tạo ra Fact mới hoặc cập nhật Fact hiện có (Forward Chaining).

---

Tiếp theo, chúng ta sẽ đi sâu vào chi tiết các lệnh định nghĩa (DDL).

\section{Định nghĩa Tri thức (KDL)}

KDL là tập hợp các lệnh dùng để định nghĩa cấu trúc, logic và các mối quan hệ trong Cơ sở tri thức (Knowledge Base). 

---

\subsection{Quản lý Cơ sở Tri thức (Knowledge Base)}

Lệnh khởi tạo hoặc xóa bỏ một vùng chứa tri thức logic.

*   \textbf{CREATE KNOWLEDGE BASE <name> [DESCRIPTION "<text>"]}: Khởi tạo KB mới.
*   \textbf{ALTER KNOWLEDGE BASE <name> SET (DESCRIPTION: "<text>")}: Cập nhật mô tả cho KB.
*   \textbf{DROP KNOWLEDGE BASE <name>}: Xóa toàn bộ dữ liệu và cấu trúc của KB.
*   \textbf{USE <name>}: Chuyển ngữ cảnh làm việc sang KB được chỉ định.

---

\subsection{Định nghĩa Khái niệm (Concept)}

Concept là thực thể cốt lõi trong KBMS, tương tự như Class trong OOP hay Table trong SQL nhưng linh hoạt hơn.

\subsubsection{Tạo Concept (CREATE CONCEPT)}
Lệnh này định nghĩa cấu trúc của một Khái niệm. Bạn có thể định nghĩa theo kiểu liệt kê biến đơn giản hoặc sử dụng các khối chuyên sâu:

\begin{verbatim}
CREATE CONCEPT <name> (
    VARIABLES (<var>: <type>, ...),
    ALIASES (<alias1>, <alias2>),
    BASE_OBJECTS (<obj1>, <obj2>),
    CONSTRAINTS (<if_then_constraints>),
    SAME_VARIABLES (<var1>, <var2>),
    CONSTRUCT_RELATIONS (<rel_definitions>),
    PROPERTIES (<prop1>, <prop2>),
    RULES (<logic_rules>),
    EQUATIONS (<math_equations>)
);
\end{verbatim}

\#\#\#\# Chi tiết các khối nâng cao:
*   \textbf{ALIASES}: Định nghĩa các tên gọi khác cho Concept.
*   \textbf{BASE\_OBJECTS}: Liệt kê các đối tượng cơ sở cấu thành nên Concept này (thường dùng trong PART\_OF).
*   \textbf{CONSTRAINTS}: Các ràng buộc logic nội bộ phải thỏa mãn.
*   \textbf{SAME\_VARIABLES}: Khai báo các biến ở các thành phần khác nhau nhưng trỏ cùng một giá trị tri thức.
*   \textbf{CONSTRUCT\_RELATIONS}: Tự động tạo lập quan hệ khi dữ liệu được nạp vào.
*   \textbf{PROPERTIES}: Các đặc tính metadata của Concept.
*   \textbf{RULES / EQUATIONS}: Định nghĩa luật và phương trình ngay bên trong Concept thay vì tạo lệnh rời.

\subsubsection{Chỉnh sửa Concept (ALTER CONCEPT)}
Hỗ trợ thay đổi cấu trúc mà không làm mất dữ liệu hiện có.

\#\#\#\# Cấu trúc lệnh
\begin{verbatim}
ALTER CONCEPT <name> (
    ADD (
        VARIABLE <var>: <type>,
        RULE (<rule_definition>),
        CONSTRAINT (<expression>),
        EQUATION '<expression>',
        PROPERTY <key>: <value>
    ),
    DROP (
        VARIABLE <name>,
        RULE <name>,
        CONSTRAINT <name>
    ),
    RENAME (VARIABLE <old> TO <new>)
);
\end{verbatim}

> [!NOTE]
> `ALTER CONCEPT` là lệnh duy nhất cho phép chỉnh sửa các thực thể nội bộ (internal rules, internal relations) mà không cần xóa đi tạo lại.

---

\subsection{Quản lý Phân cấp (Hierarchy)}

Định nghĩa quan hệ cha-con hoặc thành phần giữa các Concept.

\textit{   \textbf{IS\_A / ISA}: Quan hệ kế thừa (ví dụ: }Square IS\_A Rectangle*).
\textit{   \textbf{PART\_OF}: Quan hệ thành phần (ví dụ: }Engine PART\_OF Car*).

\subsubsection{Cú pháp}
\begin{verbatim}
ADD HIERARCHY <child> IS_A <parent>;
ADD HIERARCHY <part> PART_OF <whole>;
REMOVE HIERARCHY <parent> IS_A <child>;
\end{verbatim}

---

\subsection{Định nghĩa Quan hệ (Relation)}

Mô tả các mối liên kết ngữ nghĩa giữa hai Concept.

\subsubsection{Cú pháp}
\begin{verbatim}
CREATE RELATION <name> 
FROM <domain> TO <range>
[PARAMS (<p1>, ...)]
[PROPERTIES (<transitive, symmetric, functional, ...>)];
\end{verbatim}

---

\subsection{Định nghĩa Luật (Rule)}

Định nghĩa logic suy diễn tự động (Forward Chaining).

\subsubsection{Cú pháp}
\begin{verbatim}
CREATE RULE <name>
SCOPE <concept>
IF <condition>
THEN <action>;
\end{verbatim}
\textit{Ví dụ:} `CREATE RULE FeverCheck SCOPE Patient IF temp > 37.5 THEN SET status = 'Fever';`

---

\subsection{Toán tử & Hàm Tùy biến (Operator & Function)}

Mở rộng khả năng tính toán của hệ thống.

*   \textbf{CREATE OPERATOR <symbol>}: Định nghĩa toán tử mới (ví dụ: `|-|`).
*   \textbf{CREATE FUNCTION <name>}: Định nghĩa hàm con xử lý logic.

\subsubsection{Cú pháp chung}
\begin{verbatim}
CREATE {OPERATOR | FUNCTION} <identity>
PARAMS (<p1> <type1>, ...)
RETURNS <type>
BODY '<script_logic>';
\end{verbatim}

---

\subsection{Công thức tự động (Computation)}

Gán công thức tính toán trực tiếp vào các biến của Concept.

\subsubsection{Cú pháp}
\begin{verbatim}
ADD COMPUTATION TO <concept>
VARIABLES <inputs>, <result>
FORMULA '<expression>';
\end{verbatim}

---

\subsection{Chỉ mục (Index)}

Tối ưu hóa tốc độ truy vấn trên các biến cụ thể của Concept (sử dụng B+ Tree).

\subsubsection{Cú pháp}
\begin{verbatim}
CREATE INDEX <index_name> ON <concept_name> (<variable_list>);
DROP INDEX <index_name>;
\end{verbatim}

---

\subsection{Sự kiện tự động (Trigger)}

Tự động thực thi một câu lệnh KBQL khi có sự kiện dữ liệu xảy ra.

\subsubsection{Cú pháp}
\begin{verbatim}
CREATE TRIGGER <name>
ON ( {INSERT|UPDATE|DELETE} OF <concept> )
DO ( <kbql_statement> );
\end{verbatim}
\textit{Ví dụ:} `CREATE TRIGGER LogInsert ON (INSERT OF Patient) DO (PRINT 'New patient added');`

\section{Thao tác tri thức (KML)}

KML là tập hợp các lệnh dùng để chèn, cập nhật, xóa các Sự kiện (Facts) và chuyển đổi dữ liệu trong Cơ sở tri thức.

---

\subsection{Chèn Sự kiện (INSERT)}

Lệnh chèn một hoặc nhiều Fact vào một Concept.

\subsubsection{Chèn một dòng đơn}
\begin{verbatim}
INSERT INTO <concept_name> ATTRIBUTE (<val1>, <val2>, ...);
\end{verbatim}

\subsubsection{Chèn hàng loạt (Bulk Insert)}
Dùng để nạp dữ liệu lớn một cách nhanh chóng.
\begin{verbatim}
INSERT BULK INTO <concept_name> ATTRIBUTE (
    (<val1a>, <val2a>, ...),
    (<val1b>, <val2b>, ...),
    ...
);
\end{verbatim}

---

\subsection{Cập nhật Sự kiện (UPDATE)}

Sửa đổi giá trị các biến của các Fact hiện có dựa trên điều kiện lọc.

\subsubsection{Cú pháp}
\begin{verbatim}
UPDATE <concept_name> 
ATTRIBUTE (SET <var1>: <new_val1>, <var2>: <new_val2>) 
WHERE <conditions>;
\end{verbatim}
\textit{Lưu ý:} Khi cập nhật một thuộc tính, hệ thống sẽ tự động chạy lại các Luật (Rules) liên quan để đảm bảo tính nhất quán của tri thức (Consistency).

---

\subsection{Xóa Sự kiện (DELETE)}

Loại bỏ các Fact khỏi Concept.

\subsubsection{Cú pháp}
\begin{verbatim}
DELETE FROM <concept_name> WHERE <conditions>;
\end{verbatim}

---

\subsection{Chuyển đổi Dữ liệu (IMPORT / EXPORT)}

Công cụ sao lưu và phục hồi tri thức từ các tệp tin bên ngoài.

\subsubsection{Xuất dữ liệu (EXPORT)}
\begin{verbatim}
EXPORT (
    CONCEPT: <name>, 
    FILE: '<path>', 
    FORMAT: {CSV|JSON|XML}
);
\end{verbatim}

\subsubsection{Nhập dữ liệu (IMPORT)}
\begin{verbatim}
IMPORT (
    CONCEPT: <name>, 
    FILE: '<path>', 
    FORMAT: {CSV|JSON|XML}
);
\end{verbatim}

---

\subsection{Đặc điểm của KML trong KBMS}

Khác với các hệ quản trị CSDL thông thường, mọi thao tác KML trong KBMS đều có thể kích hoạt:
1.  \textbf{Reasoning Engine:} Các luật suy diễn tiến sẽ được kích hoạt ngay lập tức sau khi `INSERT` hoặc `UPDATE`.
2.  \textbf{Trigger Engine:} Các sự kiện tự động (`Trigger`) sẽ được thực thi nếu có định nghĩa cho sự kiện tương ứng.
3.  \textbf{Validation:} Dữ liệu được kiểm tra kiểu dữ liệu và các ràng buộc (Constraints) ngay tại tầng Parser và Storage.

\section{Truy vấn và Suy diễn (KQL)}

KQL là tập hợp các lệnh dùng để lấy thông tin, thực hiện truy vấn dữ liệu và yêu cầu hệ thống thực hiện các phép suy diễn tri thức phức tạp.

---

\subsection{Truy vấn Dữ liệu (SELECT)}

Lệnh `SELECT` trong KBQL tương tự SQL nhưng được tối ưu để làm việc với các Concept và thuộc tính tri thức.

\subsubsection{Cú pháp đầy đủ}
\begin{verbatim}
SELECT [<columns> | * | AGGREGATE(<var>)]
FROM <concept> [AS <alias>]
[JOIN <concept> ON <condition>]
[WHERE <conditions>]
[GROUP BY <vars>] 
[HAVING <condition>]
[ORDER BY <vars> {ASC|DESC}]
[LIMIT <n> OFFSET <m>];
\end{verbatim}

\subsubsection{Các tính năng đặc biệt}
*   \textbf{Hàm Tổng hợp (Aggregates):} Hỗ trợ `COUNT`, `SUM`, `AVG`, `MIN`, `MAX`.
*   \textbf{HAVING:} Mệnh đề lọc sau khi đã thực hiện `GROUP BY`.
*   \textbf{Hàm Tính toán CALC():} Thực hiện các biểu thức toán học phức tạp ngay trong lúc truy vấn.
    \textit{Ví dụ:} `SELECT name, CALC(price * 1.1) AS price\_tax FROM Product;`

---

\subsection{Giải quyết Tri thức (SOLVE)}

Đây là lệnh đặc trưng của KBMS, dùng để kích hoạt bộ máy giải quyết vấn đề (Problem Solver) dựa trên các công thức và luật đã định nghĩa.

\subsubsection{Cú pháp}
\begin{verbatim}
SOLVE ON CONCEPT <name>
GIVEN <inputs>
FIND <outputs>
[SAVE]
[USING <rule_or_relation_list>];
\end{verbatim}

\subsubsection{Cách hoạt động}
1.  \textbf{Input:} Nạp các giá trị đã biết vào phần `GIVEN`.
2.  \textbf{Inference:} Hệ thống sử dụng Forward Chaining và Equation Solving để tìm giá trị của các biến trong phần `FIND`.
3.  \textbf{Persist (Tùy chọn):} Nếu có từ khóa `SAVE`, các kết quả vừa tìm thấy sẽ được lưu trực tiếp thành một Fact mới vào Concept tương ứng.
4.  \textbf{Result:} Trả về kết quả sau khi đã suy diễn thành công.

---

\subsection{Kiểm tra Hệ thống (SHOW)}

Dùng để liệt kê các thành phần đang có trong Cơ sở Tri thức.

*   \textbf{SHOW CONCEPTS}: Liệt kê tất cả các Khái niệm.
*   \textbf{SHOW RULES}: Liệt kê các luật suy diễn.
*   \textbf{SHOW RELATIONS}: Liệt kê các quan hệ giữa các Concept.
*   \textbf{SHOW HIERARCHIES}: Xem cấu trúc phân cấp cây tri thức.
*   \textbf{SHOW OPERATORS / FUNCTIONS}: Xem các toán tử và hàm tùy biến.

---

\subsection{Phân tích & Mô tả (DESCRIBE / EXPLAIN)}

Công cụ dành cho lập trình viên để hiểu sâu về cấu trúc và cách thực thi của hệ thống.

*   \textbf{DESCRIBE {CONCEPT|KB|RULE|...} <name>}: Hiển thị chi tiết cấu trúc định nghĩa của một đối tượng.
*   \textbf{EXPLAIN (<kbql\_statement>)}: Hiển thị kế hoạch thực thi (Execution Plan) của một câu lệnh, rất hữu ích để tối ưu hóa truy vấn.

---

\subsection{Truy vấn Metadata (Hệ thống)}

Bạn có thể truy vấn trực tiếp vào danh mục hệ thống để lấy thông tin metadata.
\textit{Ví dụ:} `SELECT * FROM <concept\_name>.variables;` trả về danh sách các biến và kiểu dữ liệu của Concept đó.

\section{Kiểm soát truy cập (KCL)}

KCL là tập hợp các lệnh dùng để quản lý hệ thống bảo mật, tài khoản và quyền hạn truy cập của người dùng trong KBMS.

---

\subsection{Quản lý Người dùng (USER Management)}

KBMS cho phép định cấu hình người dùng và vai trò để bảo vệ tri thức khỏi các truy cập không mong muốn.

\subsubsection{Tạo Người dùng mới (CREATE USER)}
\begin{verbatim}
CREATE USER <username> 
PASSWORD '<password>' 
[ROLE {ADMIN|SERVICE|USER}];
\end{verbatim}

\subsubsection{Chỉnh sửa Người dùng (ALTER USER)}
Sử dụng để đổi mật khẩu hoặc trạng thái quản trị.
\begin{verbatim}
ALTER USER <username> (
    SET (PASSWORD: '<new_pass>', ADMIN: true)
);
\end{verbatim}

\subsubsection{Xóa Người dùng (DROP USER)}
\begin{verbatim}
DROP USER <username>;
\end{verbatim}

---

\subsection{Quản trị Quyền hạn (GRANT / REVOKE)}

Sử dụng để cấp hoặc thu hồi quyền thao tác trên các Concept và Tri thức cụ thể.

\subsubsection{Cấp quyền (GRANT)}
\begin{verbatim}
GRANT {SELECT, INSERT, UPDATE, DELETE, ...} 
ON CONCEPT <name> 
TO <username>;
\end{verbatim}

\subsubsection{Thu hồi quyền (REVOKE)}
\begin{verbatim}
REVOKE {SELECT, INSERT, UPDATE, DELETE, ...} 
ON CONCEPT <name> 
FROM <username>;
\end{verbatim}

---

\subsection{Các thực thể Quản trị đặc biệt (System Admin)}

*   \textbf{ADMIN}: Có toàn quyền trên tất cả các KB, Concept và người dùng khác.
*   \textbf{SERVICE}: Có quyền truy cập vào các API hệ thống nhưng không thể thay đổi cấu trúc Metadata.
*   \textbf{USER}: Quyền hạn mặc định, cần được cấp phép cụ thể cho từng đối tượng.

\section{TCL: Ngôn ngữ Kiểm soát Giao dịch (Transaction Control Language)}

TCL là tập hợp các lệnh dùng để quản lý sự thực thi đồng nhất của các lệnh KBQL, đảm bảo tính toàn vẹn tri thức (ACID).

---

\subsection{Một Giao dịch là gì? (Transaction)}

Giao dịch là một đơn vị công việc bao gồm một hoặc nhiều lệnh thực thi trên KBMS. Nếu một phần của giao dịch thất bại, toàn bộ giao dịch có thể được hủy bỏ để không làm sai lệch tri thức hiện tại.

---

\subsection{Các lệnh Giao dịch}

\subsection{Các lệnh Giao dịch}

\subsubsection{Bắt đầu Giao dịch (BEGIN TRANSACTION)}
\begin{verbatim}
BEGIN TRANSACTION;
\end{verbatim}
Từ thời điểm này, mọi thay đổi dữ liệu hoặc định nghĩa tri thức sẽ được thực thi tạm thời trong bộ đệm giao dịch.

\subsubsection{Xác nhận Giao dịch (COMMIT)}
\begin{verbatim}
COMMIT;
\end{verbatim}
Lệnh này xác nhận toàn bộ các thay đổi trong giao dịch và lưu trữ vĩnh viễn vào các file B+ Tree và Catalog của hệ thống.

\subsubsection{Hủy bỏ Giao dịch (ROLLBACK)}
\begin{verbatim}
ROLLBACK;
\end{verbatim}
Lệnh này hủy bỏ toàn bộ các lệnh từ lúc `BEGIN TRANSACTION`, đưa tri thức quay về trạng thái trước khi giao dịch bắt đầu.

---

\subsection{Vai trò của Giao dịch trong KBMS}

1.  \textbf{Tính Nguyên tử (Atomicity):} Đảm bảo một bộ các Fact liên quan được nạp vào cùng lúc (ví dụ: Thông tin bệnh nhân và các triệu chứng chẩn đoán kèm theo).
2.  \textbf{Tính Nhất quán (Consistency):} Ngăn chặn việc Rules thực hiện các thay đổi tri thức không đồng bộ giữa các Concept khác nhau.
3.  \textbf{Cách ly (Isolation):} Các giao dịch đang diễn ra không ảnh hưởng đến các truy vấn suy diễn khác cho tới khi `COMMIT` thành công.

\section{Bảo trì Hệ thống (Maintenance)}

Hệ quản trị KBMS cung cấp một bộ các lệnh quản trị chuyên sâu giúp tối ưu hóa bộ nhớ, chỉ mục và kiểm tra tính nhất quán của tri thức.

---

\subsection{Lệnh Bảo trì (MAINTENANCE)}

Lệnh `MAINTENANCE` được dùng để thực hiện các tác vụ dọn dẹp và sửa lỗi hệ thống cấp thấp.

\subsubsection{Cú pháp}
\begin{verbatim}
MAINTENANCE (
    {VACUUM | REINDEX | CHECK CONSISTENCY} [ON CONCEPT <name>],
    ...
);
\end{verbatim}

---

\subsection{Các hành động cụ thể}

\subsubsection{Thu gọn Bộ nhớ (VACUUM)}
Khi Fact bị xóa, không gian bộ nhớ trong cây B+ Tree không được giải phóng ngay lập tức. Lệnh `VACUUM` thực hiện việc sắp xếp lại các Page dữ liệu và thu hồi dung lượng đĩa cứng dư thừa.
\textit{Ví dụ:} `MAINTENANCE (VACUUM);`

\subsubsection{Đánh lại Chỉ mục (REINDEX)}
Cập nhật lại toàn bộ các Index gắn với Concept. Hữu ích khi hệ thống gặp sự cố mất điện hoặc phát hiện sai lệch chỉ mục.
\textit{Ví dụ:} `MAINTENANCE (REINDEX ON CONCEPT Patient);`

\subsubsection{Kiểm tra Tính nhất quán (CHECK CONSISTENCY)}
Kiểm tra xem các liên kết giữa Fact, Rule và Hierarchy có còn hợp lệ hay không. Đảm bảo tri thức không bị mâu thuẫn.
\textit{Ví dụ:} `MAINTENANCE (CHECK CONSISTENCY ON CONCEPT (*));`

---

\subsection{Tần suất thực hiện bản khuyến nghị}

*   \textbf{VACUUM:} Sau mỗi lần xóa dữ liệu quy mô lớn (Batch delete).
*   \textbf{REINDEX:} Định kỳ hàng tuần hoặc sau khi thực hiện việc thay đổi Metadata phức tạp.
*   \textbf{CHECK CONSISTENCY:} Thực hiện trước khi triển khai hệ thống suy diễn vào môi trường quan trọng.

\section{Biểu thức và Hàm số (Expressions)}

Biểu thức là thành phần cốt lõi xuất hiện trong các mệnh đề `WHERE`, `IF`, `SET` và hàm `CALC()`. KBQL tích hợp một bộ máy đánh giá biểu thức mạnh mẽ hỗ trợ đầy đủ các phép toán và hàm học thuật.

---

\subsection{Toán tử Cơ bản (Operators)}

\subsubsection{Toán tử Số học}
\textit{   `+`, `-`, `}`, `/`: Các phép toán cơ bản.
*   `^`: Lũy thừa.
*   `\%`: Chia lấy dư.

\subsubsection{Toán tử So sánh}
*   `=`, `!=` (hoặc `<>`): So sánh bằng và không bằng.
*   `<`, `<=`, `>`, `>=`: So sánh thứ tự.

\subsubsection{Toán tử Logic}
*   `AND`, `OR`, `NOT`: Các phép logic Boolean.

---

\subsection{Hàm Tích hợp (Built-in Functions)}

KBMS tích hợp sẵn các hàm toán học sau đây để phục vụ tính toán tri thức:

| Hàm | Giải thích | Ví dụ |
| :--- | :--- | :--- |
| `Abs(x)` | Giá trị tuyệt đối | `Abs(-10)` $\rightarrow$ 10 |
| `Sqrt(x)` | Căn bậc hai | `Sqrt(16)` $\rightarrow$ 4 |
| `Pow(x, y)` | Lũy thừa | `Pow(2, 3)` $\rightarrow$ 8 |
| `Round(x, n)` | Làm tròn đến n chữ số | `Round(3.1415, 2)` $\rightarrow$ 3.14 |
| `Floor(x)` / `Ceiling(x)` | Làm tròn xuống / lên | `Floor(2.9)` $\rightarrow$ 2 |
| `Sin`, `Cos`, `Tan` | Các hàm lượng giác | `Sin(0)` $\rightarrow$ 0 |
| `Factorial(n)` | Tính giai thừa | `Factorial(5)` $\rightarrow$ 120 |

---

\subsection{Cách sử dụng trong CALC()}

Hàm `CALC()` cho phép nhúng trực tiếp biểu thức vào kết quả trả về của `SELECT`.
\textit{Ví dụ:} `SELECT name, CALC(Sqrt(a^2 + b^2)) AS hypotenuse FROM Triangles;`

---

\subsection{Tương tác với Biến (Variables)}

Bên trong các biểu thức, bạn có thể tham chiếu trực tiếp đến tên các biến của Concept hiện tại. Nếu có `JOIN`, hãy sử dụng Alias để phân biệt: `e.salary * 1.1`.

---

\subsection{Lưu ý về Hiệu năng}

Các biểu thức phức tạp trong `WHERE` có thể làm chậm tốc độ truy vấn nếu không có Index phù hợp. Hãy sử dụng lệnh `EXPLAIN` để kiểm tra cách hệ thống đánh giá các biểu thức của bạn.

\section{Các trường hợp Chuyên gia (Expert Cases)}

Tài liệu này trình bày các tình huống sử dụng KBQL ở mức độ nâng cao, khai thác tối đa sức mạnh của bộ máy suy diễn và giải toán của KBMS.

---

\subsection{Giải ngược Phương trình (Inverse Problem Solving)}

Hệ thống KBMS không chỉ tính toán một chiều mà còn có khả năng giải ngược các tham số đầu vào khi biết kết quả đầu ra thông qua các thuật toán \textbf{Newton-Raphson} và \textbf{Brent}.

\subsubsection{Kịch bản: Tính toán mạch điện}
Chúng ta biết luật Ohm `U = I * R`. Nếu biết `U` và `I`, hệ thống phải tự tìm `R`.

\begin{verbatim}
CREATE CONCEPT Resistor (
    VARIABLES (u: DECIMAL, i: DECIMAL, r: DECIMAL),
    EQUATIONS (u = i * r)
);

-- Trường hợp 1: Biết I và R, tìm U (Tính xuôi)
SOLVE ON CONCEPT Resistor GIVEN i = 2, r = 10 FIND u; -- Kết quả: u = 20

-- Trường hợp 2: Biết U và I, tìm R (Tính ngược - Expert)
SOLVE ON CONCEPT Resistor GIVEN u = 220, i = 2 FIND r; -- Kết quả: r = 110
\end{verbatim}

---

\subsection{Suy diễn Chùm (Chained Inference)}

Thể hiện khả năng lan truyền tri thức qua nhiều lớp luật (Forward Chaining).

\subsubsection{Kịch bản: Hệ thống Cảnh báo sớm}
\begin{verbatim}
CREATE RULE R1 SCOPE Sensor IF temp > 100 THEN SET status = 'Overheat';
CREATE RULE R2 SCOPE Sensor IF status = 'Overheat' AND pressure > 50 THEN SET alert = 'Critical';
CREATE RULE R3 SCOPE Sensor IF alert = 'Critical' THEN DO (PRINT 'EMERGENCY SHUTDOWN');

-- Khi chèn dữ liệu thỏa mãn chuỗi:
INSERT INTO Sensor ATTRIBUTE (110, 60);
-- Kết quả: R1 kích hoạt -> status='Overheat' -> R2 kích hoạt -> alert='Critical' -> R3 kích hoạt -> In ra thông báo.
\end{verbatim}

---

\subsection{Phân cấp Kế thừa Đa tầng (Recursive Property Propagation)}

Khi định nghĩa Hierarchy, các luật ở Concept cha sẽ tự động "chảy" xuống toàn bộ các con, cháu.

\subsubsection{Kịch bản: Phân loại Hình học}
\begin{verbatim}
CREATE CONCEPT Shape (VARIABLES (color: STRING));
CREATE CONCEPT Polygon (...);
CREATE CONCEPT Triangle (...);

ADD HIERARCHY Polygon IS_A Shape;
ADD HIERARCHY Triangle IS_A Polygon;

CREATE RULE ShapeColor SCOPE Shape IF color = 'Red' THEN SET is_hot = true;

-- Khi chèn một Triangle màu đỏ:
INSERT INTO Triangle ATTRIBUTE ('Red', ...);
-- Kết quả: Mặc dù luật định nghĩa ở Shape, nhưng Triangle vẫn nhận được is_hot = true.
\end{verbatim}

---

\subsection{Ràng buộc nội bộ phức tạp (Internal Constraints)}

Sử dụng khối `CONSTRAINTS` để bảo vệ dữ liệu ở mức tri thức, không cho phép các Fact mâu thuẫn tồn tại.

\subsubsection{Kịch bản: Quản lý nhân sự}
\begin{verbatim}
CREATE CONCEPT Employee (
    VARIABLES (age: INT, experience: INT),
    CONSTRAINTS (
        IF age < experience + 18 THEN DO (ERROR 'Tuổi không hợp lệ so với kinh nghiệm')
    )
);

-- Lệnh này sẽ bị Parser/Engine chặn lại:
INSERT INTO Employee ATTRIBUTE (20, 5); -- (20 < 5 + 18) => Lỗi.
\end{verbatim}

---

\subsection{Giải hệ phương trình 2 ẩn (Newton-Raphson 2D)}

KBMS hỗ trợ giải đồng thời các hệ thức toán học.

\subsubsection{Kịch bản: Tìm giao điểm đường thẳng}
\begin{verbatim}
CREATE CONCEPT Intersection (
    VARIABLES (x: DECIMAL, y: DECIMAL),
    EQUATIONS (
        y = 2 * x + 1,
        y = -x + 4
    )
);

SOLVE ON CONCEPT Intersection FIND x, y;
-- Kết quả: x = 1, y = 3 (Hệ thống tự giải hệ phương trình).
\end{verbatim}

---

\subsection{Kết hợp Trigger và Rule (The Chain Reaction)}

Dùng Trigger để kích hoạt một luồng suy diễn ở một Concept hoàn toàn khác.

\subsubsection{Kịch bản: Đồng bộ kho và đơn hàng}
\begin{verbatim}
CREATE TRIGGER SyncStock ON (INSERT OF Order) 
DO (UPDATE Product ATTRIBUTE (SET stock: stock - 1) WHERE id = new.product_id);

-- Khi cập nhật kho, các luật về 'Out of stock' sẽ tự động chạy:
CREATE RULE OutOfStock SCOPE Product IF stock < 0 THEN SET status = 'Refill';
\end{verbatim}

\section{Xác thực Ngôn ngữ (Language Validation)}

Để đảm bảo mọi câu lệnh KBQL đều được diễn dịch chính xác trước khi gửi tới Engine, hệ thống KBMS trải qua quy trình kiểm thử cú pháp nghiêm ngặt.

\subsection{Kiểm thử Cú pháp (Parsing Test)}

Tầng Parser được kiểm soát bởi tệp \textbf{`ParserTests.cs`} với hơn 21,000 dòng kịch bản kiểm thử, bao phủ mọi trường hợp từ câu lệnh đơn giản đến phức tạp nhất.

| Nhóm Kiểm thử | Mô tả | Trạng thái |
| :--- | :--- | :--- |
| \textbf{KDL Validation} | Kiểm tra định nghĩa Concept, Hierarchy, Relation. | \textbf{Passed} |
| \textbf{KQL Validation} | Kiểm tra các phép toán Join, Filter, Order By, Limit. | \textbf{Passed} |
| \textbf{Expression Eval} | Kiểm tra các biểu thức toán học phức tạp (`CALC()`). | \textbf{Passed} |
| \textbf{Error Handling} | Kiểm tra khả năng báo lỗi đúng dòng, đúng cột. | \textbf{Passed} |

\subsection{Bằng chứng Thực nghiệm (Evidence)}

Kết quả chạy bộ kiểm thử `ParserTests`:

\textbf{[Missing Image: placeholder_parser_test_log.png]}\\ \textit{Hình 6.1: Kết quả kiểm thử bộ phân tích cú pháp (Parser) với hơn 21,000 kịch bản.}

\subsection{Nhật ký Lexer (Lexing Proof)}

Trước khi phân tích cú pháp, Lexer thực hiện tách từ tố (Tokenization). Tệp `LexerTests.cs` đảm bảo mọi ký tự đặc biệt và từ khóa đều được nhận diện đúng.

\textbf{[Missing Image: placeholder_lexer_token_stream.png]}\\ \textit{Hình 6.2: Luồng Token được phân tách bởi Lexer trước khi đưa vào Parser.}

---

> [!NOTE]
> Việc vượt qua bộ kiểm thử `ParserTests` đảm bảo rằng KBMS sẽ không bao giờ thực thi một câu lệnh sai cú pháp, giúp bảo vệ an toàn cho tầng lưu trữ.

